\documentclass[12pt]{article}
\usepackage{amsmath, amssymb}
\usepackage{enumitem}
\usepackage{geometry}
\geometry{margin=1in}

\begin{document}

\section*{Q1. Division into Equal Groups (Multinomial Counting)}

\subsection*{Problem Structure}

\begin{itemize}
    \item 20 distinct dishes.
    \item Divide into 4 groups of 5 dishes each.
    \item Order within groups does not matter.
    \item Order of groups does not matter.
\end{itemize}

\subsection*{(a) Total Number of Ways}

Since:

\begin{itemize}
    \item All 20 dishes are distinct.
    \item Groups are unlabeled.
    \item Each group has 5 elements.
\end{itemize}

We apply the \textbf{multinomial coefficient}:

\[
\text{Total ways} = \frac{20!}{(5!)^4 4!}
\]

\textbf{Reasoning:}

\begin{itemize}
    \item (20!) counts all permutations.
    \item Divide by (5!) for each group (order within group irrelevant).
    \item Divide by (4!) because groups themselves are unlabeled.
\end{itemize}

\subsection*{(b) Probability All Platters Have Same Cuisine}

Given:

\begin{itemize}
    \item 5 dishes from each cuisine.
    \item Each cuisine must form exactly one platter.
\end{itemize}

Since platters are unlabeled, there is \textbf{exactly one favorable arrangement}.

\[
P = \frac{1}{\frac{20!}{(5!)^4 4!}} = \frac{(5!)^4 4!}{20!}
\]

\section*{Q2. Uniform Distribution \& Waiting Time}

Arrival time $\sim$ Uniform over 30 minutes (7:00--7:30).

\subsection*{Key Idea:}

Probability = Favorable time / Total time

\subsection*{(a) Wait $<$ 5 minutes}

Arrival must be in:

\begin{itemize}
    \item 7:10--7:15
    \item 7:25--7:30
\end{itemize}

Total favorable time = 10 minutes.

\[
P(\text{wait} < 5) = \frac{10}{30} = \frac{1}{3}
\]

\subsection*{(b) Wait $>$ 10 minutes}

Arrival must be:

\begin{itemize}
    \item 7:00--7:05
    \item 7:15--7:20
\end{itemize}

Total favorable time = 10 minutes.

\[
P(\text{wait} > 10) = \frac{10}{30} = \frac{1}{3}
\]

\section*{Q3. Conditional Probability with Complements}

Let:

\begin{itemize}
    \item (L) = arrives late
    \item (E) = leaves early
\end{itemize}

Given:

\[
P(L)=0.15,\quad P(E)=0.25,\quad P(L\cap E)=0.08
\]

We want:

\[
P(L^c \mid E^c)
\]

\subsection*{Step 1: Complements}

\[
P(E^c)=1-0.25=0.75
\]

\subsection*{Step 2: Use Inclusion--Exclusion}

\[
P(L\cup E)=0.15+0.25-0.08=0.32
\]

\subsection*{Step 3: De Morgan’s Law}

\[
P(L^c \cap E^c)=1-P(L\cup E)=1-0.32=0.68
\]

\subsection*{Step 4: Conditional Probability}

\[
P(L^c|E^c)=\frac{0.68}{0.75}=0.907
\]

\section*{Q4. Poisson Distribution ($\lambda = 0.2$)}

Poisson PMF:

\[
P(X=k)=e^{-\lambda}\frac{\lambda^k}{k!}
\]

\subsection*{(a) $P(X=0)$}

\[
P(X=0)=e^{-0.2}=0.8187
\]

\subsection*{(b) $P(X\ge2)$}

\[
P(X\ge2)=1-P(0)-P(1)
\]

\[
P(1)=0.2e^{-0.2}=0.1637
\]

\[
P(X\ge2)=1-0.8187-0.1637=0.0176
\]

\section*{Q5. Poisson Distribution ($\lambda = 3.5$)}

\subsection*{(a) $P(X\ge2)$}

\[
P(X\ge2)=1-P(0)-P(1)
\]

\[
P(0)=e^{-3.5}=0.0302
\]

\[
P(1)=3.5e^{-3.5}=0.1057
\]

\[
P(X\ge2)=1-0.0302-0.1057=0.8641
\]

\subsection*{(b) $P(X\le1)$}

\[
P(X\le1)=0.0302+0.1057=0.1359
\]

\section*{Q6. Poisson ($\lambda = 3$)}

\subsection*{(a) $P(X\ge3)$}

\[
P(X\ge3)=1-P(0)-P(1)-P(2)
\]

\[
P(0)=0.0498
\]

\[
P(1)=0.1494
\]

\[
P(2)=0.2240
\]

\[
P(X\ge3)=1-0.4232=0.5768
\]

\subsection*{(b) Conditional Probability}

\[
P(X\ge3|X\ge1)=\frac{P(X\ge3)}{P(X\ge1)}
\]

\[
P(X\ge1)=1-P(0)=0.9502
\]

\[
= \frac{0.5768}{0.9502}=0.607
\]

\section*{Q7. Gaussian Distribution}

Given PDF:

\[
f_X(x)=\frac{1}{\sqrt{8\pi}} e^{-(x+3)^2/8}
\]

Compare with standard Gaussian:

\[
m=-3,\quad \sigma=2
\]

\subsection*{Gaussian Relations}

\[
F_X(x)=\Phi\left(\frac{x-m}{\sigma}\right)
\]

\[
P(X>x)=Q\left(\frac{x-m}{\sigma}\right)
\]

\subsection*{Results}

\begin{enumerate}
    \item $P(X\le0)=\Phi(3/2)=1-Q(3/2)$
    \item $P(X>4)=Q(7/2)$
    \item $|X+3|<2\Rightarrow -5<X<-1$
    
    \[
    = \Phi(1)-\Phi(-1)
    \]
    
    Using symmetry:
    
    \[
    =2\Phi(1)-1=1-2Q(1)
    \]
    
    \item $|X-2|>1\Rightarrow X<1 \text{ or } X>3$
    
    \[
    =F(1)+P(X>3)=1-Q(2)+Q(3)
    \]
\end{enumerate}

\section*{Q8. Jury Decision (Conditioning)}

Let:

\begin{itemize}
    \item Each juror correct with probability ( $\theta$ )
    \item Defendant guilty with probability ( $\alpha$ )
\end{itemize}

If guilty:

\[
\sum_{i=8}^{12}\binom{12}{i}\theta^i(1-\theta)^{12-i}
\]

If innocent:

\[
\sum_{i=5}^{12}\binom{12}{i}\theta^i(1-\theta)^{12-i}
\]

Total probability:

\[
\alpha \sum_{i=8}^{12} \binom{12}{i}\theta^i(1-\theta)^{12-i}
+
(1-\alpha)\sum_{i=5}^{12} \binom{12}{i}\theta^i(1-\theta)^{12-i}
\]

\section*{Q9. Determining Constant ($c$)}

Given:

\[
p(i)=c\frac{\lambda^i}{i!}
\]

Using:

\[
\sum_{i=0}^{\infty}p(i)=1
\]

Since:

\[
e^\lambda=\sum_{i=0}^{\infty}\frac{\lambda^i}{i!}
\]

Thus:

\[
ce^\lambda=1
\]

\[
c=e^{-\lambda}
\]

\subsection*{Results}

\[
P(X=0)=e^{-\lambda}
\]

\[
P(X>2)=1-e^{-\lambda}-\lambda e^{-\lambda}-\frac{\lambda^2}{2}e^{-\lambda}
\]

\subsection*{CDF for Discrete Variable}

\[
F(a)=\sum_{x\le a}p(x)
\]

It is a \textbf{step function} with jumps equal to (p($x_i$)).

\section*{Q10. System Reliability (Binomial)}

Each component functions with probability (p).

\subsection*{(a) Compare 5 vs 3 Components}

5-component effective if $\ge3$ work:

\[
10p^3(1-p)^2+5p^4(1-p)+p^5
\]

3-component effective if $\ge2$ work:

\[
3p^2(1-p)+p^3
\]

Difference reduces to:

\[
3(p-1)^2(2p-1)>0
\]

Thus:

\[
p>\frac{1}{2}
\]

\subsection*{(b) General Case}

Difference:

\[
\binom{2k-1}{k} p^k(1-p)^k (2p-1)
\]

Since all other terms positive:

\[
> 0 \iff p>\frac{1}{2}
\]

\end{document}